\chapter{Droits, provenance et retrait}\label{ch:droits}

Un projet qui travaille sur des documents d'archives, des manuscrits ou des textes
d'autres auteurs manie des droits qu'il ne détient pas. Le dépôt public rend chaque erreur
visible et durable. Ce chapitre décrit comment séparer ce qui peut être publié de ce qui
doit rester privé, et comment garder la possibilité de retirer.

\section{Trois niveaux}

\begin{description}
  \item[Public] le code, les outils, les données dérivées qu'on a le droit de diffuser,
        les articles ;
  \item[Privé] les copies locales de documents, les transcriptions de textes protégés,
        tant que les ayants droit n'ont pas donné leur accord ;
  \item[Jamais stocké] les images des documents eux-mêmes, quand la source les sert déjà :
        on les affiche depuis la source, on ne les recopie pas.
\end{description}

\section{Le fichier \texttt{.gitignore}}

Un dossier inscrit dans \fichier{.gitignore} n'entre jamais dans le dépôt : \texttt{git}
l'ignore, et rien de ce qu'il contient ne peut être poussé par mégarde. Les transcriptions
privées et les copies locales y vont.

\begin{warningblock}
\cmd{git add -f} passe outre le \fichier{.gitignore}. On ne l'utilise jamais sur un dossier
privé. Un fichier poussé une fois reste dans l'histoire même si on le supprime ensuite.
\end{warningblock}

\section{Un interrupteur à la compilation}

Quand la question des droits n'est pas réglée, on peut garder les documents dans le dépôt
sans les montrer sur le site. Une variable lue au moment de la construction du site (par
exemple \texttt{RENDER\_TRANSCRIPTS=false}) retire l'affichage en un réglage, sans toucher
au contenu. Si les ayants droit le demandent, on retire ; s'ils l'accordent, on rallume.

\section{Citations courtes}

Dans les tickets, les messages de commit et les articles, on cite quelques mots à la fois,
avec la source précise, et l'on résume le reste dans ses propres mots. La règle vaut pour
l'assistant comme pour soi, et on l'écrit dans ses procédures.

\section{Ne pas désigner les personnes}

Les dépôts publics sont indexés. On n'y nomme pas les personnes privées qui apparaissent
dans les documents, ni les correspondants qui ont répondu à une demande sans s'attendre à
être cités. On décrit leur rôle, pas leur nom.
